\chapter{Conclusion and Summary}
\thispagestyle{plain}

In conclusion, the aim of this thesis was to fabricate and characterize organic field-effect
transistors in a geometry suitable for contact-independent transport analysis.
The motivation was that conventional two-terminal OFET measurements do not probe
only the semiconductor channel. Contact voltage drops, gate-dependent injection,
trap-related relaxation, and possible leakage can all change the apparent slope
of transfer curves and therefore distort extracted mobility and threshold
voltage values~\cite{Natali2012ChargeInjection,Bittle2016MobilityOverestimation}.
For this reason, the gated van der Pauw method was used as the main analysis
tool, because it determines the sheet conductance from voltage probes that do
not carry the imposed current~\cite{Rolin2017GVDP,Lee2024GVDPIGZO}.

The theoretical discussion established the semiconductor and device-physics
background needed for this comparison. It showed why ideal FET relations are
only valid when contact effects, localized states, and leakage are sufficiently
small, and introduced the gated van der Pauw method as a four-terminal
alternative for measuring the accumulated sheet conductance. The fabricated
devices used a bottom-gate architecture with a Cr gate electrode, a PaN/PS gate
dielectric stack, a C8-BTBT molecular semiconductor layer, and MoO$_x$/Au
injection contacts. The clover-leaf contact geometry made it possible to perform
both conventional source--drain measurements and gated van der Pauw measurements
on the same device type.

The two-terminal current--voltage measurements confirmed that the fabricated
devices operate as $p$-type OFETs. Negative gate bias accumulated holes at the
semiconductor--dielectric interface, the gate leakage remained small compared
with the source--drain current, and the output curves showed the expected
ordering with gate voltage. At the same time, the gradual turn-on of the
transfer curves and the finite background resistance indicated that contact and
injection effects remained important. Thus, the conventional measurements
verified transistor action, but also showed that two-terminal mobility
extraction alone is not sufficient for this device system.

The gated van der Pauw measurements provided the main quantitative result. Sheet
conductance curves measured at $I_{34}=\SI{-100}{\nano\ampere}$,
$\SI{-200}{\nano\ampere}$, and $\SI{-300}{\nano\ampere}$ collapsed onto a common
linear dependence on effective gate voltage. The extracted mobilities were
\SI{3.11 \pm 0.02}{\centi\metre\squared\per\volt\per\second},
\SI{3.11 \pm 0.02}{\centi\metre\squared\per\volt\per\second}, and
\SI{3.10 \pm 0.02}{\centi\metre\squared\per\volt\per\second}, with threshold
voltages clustered around \SI{-10.2}{\volt} and $R^2>0.998$ for all fits. The
agreement between current levels shows that the extracted sheet conductance is
not an artefact of the chosen drive current. The most reliable mobility obtained
in this work is therefore approximately
\SI{3.1}{\centi\metre\squared\per\volt\per\second}.

The additional electrical measurements complemented this static analysis. The
DLTS-type experiment showed a fast capacitive response followed by slower
current relaxation, consistent with localized states that exchange carriers with
the channel over a distribution of time scales~\cite{Lang1974DLTS,
Hirwa2015TrapStates,Bobbert2012OperationalStability}. The effective capacitance
remained close to \SI{1.7}{\nano\farad}, while the relaxation tail indicated
that traps and bias history can affect transfer characteristics. The amplifier
test demonstrated inverted voltage amplification, but also showed that the
dynamic response is limited by large circuit resistance, parasitic capacitance,
and trap-related relaxation.

Overall, the fabricated C8-BTBT OFETs showed clear field-effect behaviour, but
their two-terminal characteristics contained substantial contact and
time-dependent contributions. The gated van der Pauw method separated the
sheet-conductance measurement from the strongest contact voltage losses and
yielded a stable, current-independent mobility. Within the scope of this thesis,
the main objective was therefore achieved: a working OFET platform was
fabricated, its nonideal two-terminal behaviour was identified, and a robust
contact-independent mobility was extracted.

Several limitations remain. The number of devices and current levels was
limited, so a broader statistical study is needed to separate device-to-device
variation from systematic material behaviour. The numerical simulation used a
simplified two-dimensional model and did not include microscopic injection
physics, trap energetics, or the full gate-stack electrostatics. Future work
should therefore combine gated van der Pauw measurements with temperature,
contact-metal, sweep-rate, and device-statistics studies, and with more complete
drift-diffusion--Poisson simulations that include contact injection and
trap-state distributions.